The current LiionTR thermal formulation uses a zero-dimensional
lumped-capacitance model. The complete battery cell is represented by a
single temperature state, such that internal spatial temperature
gradients are neglected.

The thermal energy balance is

\begin{equation}
    C_{\mathrm{th}}
    \frac{dT}{dt}
    =
    \dot{Q}_{\mathrm{gen}}
    -
    h A
    \left(
        T - T_{\infty}
    \right),
    \label{eq:lumped_thermal_balance}
\end{equation}

where $C_{\mathrm{th}}$ is the cell thermal capacity,
$\dot{Q}_{\mathrm{gen}}$ is the total internal heat-generation rate,
$h$ is the convective heat-transfer coefficient, $A$ is the external
cell surface area, and $T_{\infty}$ is the ambient temperature.

The resulting temperature derivative is

\begin{equation}
    \frac{dT}{dt}
    =
    \frac{
        \dot{Q}_{\mathrm{gen}}
        -
        h A
        \left(
            T-T_{\infty}
        \right)
    }{
        C_{\mathrm{th}}
    }.
    \label{eq:lumped_temperature_derivative}
\end{equation}


\subsection{Cell thermal capacity}

The effective thermal capacity is calculated as

\begin{equation}
    C_{\mathrm{th}}
    =
    m_{\mathrm{cell}}
    c_p,
    \label{eq:thermal_capacity}
\end{equation}

where $c_p$ is the effective specific heat capacity of the cell.

The cell mass is calculated from the effective density and geometry,

\begin{equation}
    m_{\mathrm{cell}}
    =
    \rho V,
    \label{eq:cell_mass}
\end{equation}

leading to

\begin{equation}
    C_{\mathrm{th}}
    =
    \rho V c_p.
\end{equation}

In the present implementation, density and specific heat capacity are
evaluated at \SI{298.15}{\kelvin} when cell mass and thermal capacity
are constructed. These quantities are therefore constant during a
simulation.


\subsection{Convective heat transfer}

Heat exchange with the environment is currently represented using
Newton's law of cooling,

\begin{equation}
    \dot{Q}_{\mathrm{conv}}
    =
    h A
    \left(
        T - T_{\infty}
    \right).
    \label{eq:convective_heat_loss}
\end{equation}

When $T>T_{\infty}$, this term removes heat from the cell.

If $T<T_{\infty}$, the same expression becomes negative and therefore
represents heating of the cell by the surrounding environment.


\subsection{Internal heat generation}

Internal heat generation is supplied to the thermal model by the
configured chemistry backend.

For a reaction network containing $N$ thermal runaway reactions,

\begin{equation}
    \dot{Q}_{\mathrm{gen}}
    =
    m_{\mathrm{cell}}
    \sum_{i=1}^{N}
    w_i
    \Delta H_i
    \dot{\alpha}_i.
    \label{eq:thermal_reaction_heat}
\end{equation}

The thermal model is therefore independent from the detailed kinetic
mechanism. Different reaction networks can be coupled to the same
thermal formulation.


\subsection{Thermal runaway feedback}

The coupling between reaction kinetics and temperature introduces a
positive feedback mechanism.

For Arrhenius kinetics,

\begin{equation}
    k(T)
    =
    A
    \exp
    \left(
        -\frac{E_a}{RT}
    \right).
\end{equation}

As temperature rises, the reaction rate generally increases. This
increases heat generation, which produces further temperature increase.

Conceptually,

\begin{equation}
    T
    \uparrow
    \quad \Longrightarrow \quad
    k(T)
    \uparrow
    \quad \Longrightarrow \quad
    \dot{Q}_{\mathrm{gen}}
    \uparrow
    \quad \Longrightarrow \quad
    T
    \uparrow.
\end{equation}

This feedback constitutes the principal mechanism of thermal runaway in
the present reduced-order framework.


\subsection{Validity of the lumped approximation}

The lumped-capacitance approximation assumes that the internal
temperature field equilibrates sufficiently rapidly compared with heat
transfer across the external boundary.

A conventional criterion is the Biot number,

\begin{equation}
    \mathrm{Bi}
    =
    \frac{
        h L_c
    }{
        k
    },
\end{equation}

where $L_c$ is a characteristic cell length and $k$ is the effective
thermal conductivity.

The lumped approximation is generally most appropriate when

\begin{equation}
    \mathrm{Bi}
    \ll
    1.
\end{equation}

LiionTR does not currently enforce a Biot-number criterion. Consequently,
the validity of the lumped approximation must be assessed for the
geometry, boundary conditions, and thermal runaway regime being
simulated.


\subsection{Current assumptions}

The present thermal formulation assumes:

\begin{itemize}
    \item a spatially uniform cell temperature;
    \item constant effective thermophysical properties;
    \item constant ambient temperature;
    \item uniform convection over the complete external cell surface;
    \item no radiative heat transfer;
    \item no explicit conductive coupling to neighbouring cells or
          structures;
    \item no explicit phase-change enthalpy;
    \item no internal thermal gradients.
\end{itemize}


\subsection{Limitations and future extensions}

A single-temperature model cannot reproduce internal hot spots or
temperature gradients that may develop during severe thermal runaway.

Future thermal-model developments may therefore include:

\begin{itemize}
    \item temperature-dependent density, heat capacity, and thermal
          conductivity;
    \item radiative heat transfer;
    \item conductive coupling between cells;
    \item one-dimensional or multidimensional thermal models;
    \item cell-to-cell thermal propagation;
    \item latent-heat and phase-change effects;
    \item energy coupling to vented material and gas enthalpy.
\end{itemize}